\section{Conclusiones y recomendaciones de negocio}
El proyecto recorrió en su totalidad el ciclo CRISP-DM definido en la Entrega 1, partiendo del entendimiento del problema educativo en Colombia y llegando hasta la construcción y evaluación de modelos predictivos y de segmentación sobre un cluster Big Data propio. Esta sección consolida los hallazgos del análisis exploratorio y de los modelos, y los traduce en recomendaciones concretas para el Ministerio de Educación Nacional, cerrando así el ciclo analítico iniciado en la primera entrega.

\subsection{Síntesis de los hallazgos exploratorios}
El análisis exploratorio sobre la vista minable confirmó cuantitativamente las hipótesis planteadas en la Entrega 1 y aportó matices importantes. La correlación entre el porcentaje de estudiantes con internet en el hogar y el puntaje global promedio del municipio es de $r = +0.41$, una relación moderada pero clara que se traduce en una diferencia de aproximadamente 22 puntos entre municipios del cuartil superior y del cuartil inferior de conectividad (Q1 y Q4). Esta brecha es estadísticamente robusta y consistente con la literatura sobre brecha digital y rendimiento académico citada en la Entrega 1.

Sin embargo, el hallazgo más relevante del análisis exploratorio es que la pobreza multidimensional incide en el rendimiento académico \textbf{de manera independiente al nivel de conectividad}. Al estratificar los municipios por cuartil de internet y analizar dentro de cada estrato la correlación entre el índice de privación Sisbén y el puntaje, todas las correlaciones se mantienen negativas y significativas, con valores entre $-0.40$ y $-0.52$. Esto demuestra que la conectividad no es un sustituto de las políticas sociales: incluso entre municipios con buena conectividad, la pobreza sigue penalizando el rendimiento.

El análisis territorial agregó una capa adicional. En el cuartil superior de conectividad, los municipios rurales alcanzan un puntaje promedio (257.59) que iguala e incluso supera ligeramente al promedio urbano (251.21). Esto sugiere que la conectividad rural de alto nivel tiene un efecto compensatorio particularmente fuerte, y constituye un argumento poderoso para concentrar la inversión digital en zonas rurales en lugar de seguir reforzando capitales y centros urbanos donde el techo del beneficio marginal es menor. El análisis individual a nivel estudiante reforzó esta lectura: el acceso a computador en el hogar mejora el puntaje en todas las áreas evaluadas, con un efecto particularmente fuerte en inglés ($+8.25$ puntos), lo que sugiere que el computador no solo facilita el acceso a contenidos, sino que habilita prácticas autónomas específicas.

\subsection{Síntesis de los resultados de modelado}
Los modelos supervisados convergieron en una conclusión consistente: las dos variables más explicativas del puntaje son el índice de privación Sisbén (con importancia relativa del 48.3\% en el árbol y coeficiente $-8.76$ en el modelo lineal) y la penetración de internet en el hogar (importancia del 23.3\% y coeficiente $+5.58$). La regresión lineal alcanza un $R^2$ de $0.41$ en test, mientras que el árbol de decisión llega a $0.52$, mostrando que existen interacciones no lineales que el modelo lineal no captura pero que el árbol sí logra modelar.

El modelo K-Means con $k = 2$ produjo una segmentación natural muy útil desde el punto de vista de política pública. El primer cluster (603 municipios) agrupa a los territorios de alta vulnerabilidad: bajo puntaje (227.6), baja conectividad (38.6\%) y alta privación (4.39 sobre 15). El segundo cluster (495 municipios) contiene a los territorios de baja vulnerabilidad: puntaje superior (251.4), mejor conectividad (64.2\%) y menor privación (3.23). La diferencia de 24 puntos entre los dos grupos confirma, desde el ángulo no supervisado, la misma brecha que detectaron el análisis exploratorio y los modelos supervisados.

\subsection{Recomendaciones para el Ministerio de Educación Nacional}
A partir del análisis integrado, se proponen cuatro recomendaciones de política pública orientadas a transformar los hallazgos cuantitativos en líneas de acción concretas.

\textit{Primera recomendación: focalizar la inversión en el Cluster 0.} Los 603 municipios del Cluster 0 constituyen el grupo objetivo prioritario. La estrategia debe combinar inversión en conectividad de calidad (no solo cobertura básica, sino velocidad y estabilidad) con intervenciones sociales que aborden las dimensiones del índice de privación que más pesan (educación, condiciones de la vivienda, acceso a servicios). El proyecto demuestra que invertir solo en conectividad, sin trabajar el componente socioeconómico, produce resultados subóptimos: los 10 municipios anómalos identificados en la pregunta Q8 son evidencia directa de este fenómeno.

\textit{Segunda recomendación: priorizar la conectividad rural de alto nivel.} El hallazgo de que los municipios rurales con conectividad en el cuartil superior alcanzan puntajes equivalentes a los urbanos del mismo cuartil sugiere que la inversión en zonas rurales tiene un retorno marginal alto y aún no agotado. Las políticas públicas deberían establecer metas explícitas de penetración mayor al 50\% en municipios rurales del Cluster 0, acompañadas de programas de alfabetización digital tanto para estudiantes como para acompañantes en el hogar.

\textit{Tercera recomendación: dotación de computadores con apoyo pedagógico.} El efecto del computador en el hogar es positivo en todas las áreas evaluadas y se acentúa en inglés, materia que requiere práctica autónoma con recursos multimedia. Programas tipo "Computadores para Educar" que se complementen con acompañamiento pedagógico (no solo entrega de equipos) y con licencias de contenido especializado en lengua extranjera pueden capitalizar este efecto comprobado.

\textit{Cuarta recomendación: sistema de alerta temprana basado en el modelo predictivo.} El modelo de árbol de decisión entrenado permite, en su forma actual, predecir con un error de aproximadamente $\pm 15$ puntos el puntaje promedio de un municipio dado su perfil contextual. Esta capacidad puede operativizarse en un tablero de gestión territorial que alimente al Ministerio con proyecciones para el próximo periodo Saber 11, permitiendo focalizar recursos antes de que se materialicen los resultados.

\subsection{Limitaciones y trabajo futuro}
El análisis presenta algunas limitaciones que conviene reconocer explícitamente. Primero, la unidad de análisis principal es municipio-año, lo que oculta heterogeneidad intra-municipal: dentro de un mismo municipio pueden coexistir realidades muy distintas entre colegios oficiales y privados, entre zonas urbanas y veredas, etc. Segundo, el dataset MEN tiene una variable de cobertura informativa pero no captura la calidad educativa propiamente dicha (ratio docente-estudiante, formación docente, infraestructura escolar específica), por lo que el modelo subestima el peso de factores cualitativos de la oferta educativa. Tercero, el análisis se limita al periodo cubierto por las fuentes disponibles, sin considerar choques estructurales como el cierre de instituciones durante la pandemia de COVID-19, cuyo impacto persistente sobre los resultados Saber 11 está documentado en la literatura.

Como trabajo futuro, el equipo identifica tres líneas de extensión. La primera es la incorporación de datos de calidad docente y de infraestructura escolar para enriquecer la dimensión educativa del modelo. La segunda es la construcción de un modelo a nivel estudiante individual que permita predecir resultados antes de la prueba con suficiente granularidad para alertar a colegios específicos. Esta extensión se exploró parcialmente en el bono de aprendizaje profundo de la presente entrega. La tercera línea es la integración de variables longitudinales que permitan analizar la trayectoria de cada municipio en el tiempo, evaluando el efecto causal (no solo correlacional) de las políticas implementadas.

En síntesis, el proyecto cumplió con los objetivos de la Entrega 2 al completar el ciclo CRISP-DM sobre un cluster Big Data propio, responder las ocho preguntas de negocio planteadas en la Entrega 1, entrenar tres modelos de aprendizaje de máquina con sus respectivas pruebas de hiperparámetros, y producir recomendaciones de política pública sustentadas en evidencia cuantitativa. El pipeline construido es reproducible, escalable y constituye una base sólida para futuras iteraciones del proyecto.
